\providecommand\tightlist{\setlength{\itemsep}{0pt}\setlength{\parskip}{0pt}}
\section*{1. Stance}

\begin{enumerate}

\tightlist
\item
  \textit{Go to the first-degree source.} Edit from the manuscript, print or facsimile, never from a modern edition when the source can be seen.
\item
  \textit{Distrust critical editions, this one included.} A modern edition is a witness to its editor's choices, not to the composer's.
\item
  \textit{A modern edition is not evidence of a composer's habits} below the level of structure. Spijker's editions of Wacław were reliable for repeat marks and melisma placement at cadences, and unreliable for syllable-by-syllable placement. Underlay statistics must come from period prints or manuscripts.
\item
  \textit{The music decides when the evidence runs out.} Underlay in particular must be ‘above all musical and meaningful’; if that means departing from the source, depart and say so.
\item
  \textit{Say what is editorial.} Every category of editorial intervention is stated in the method section and signalled in the score (section 8).
\end{enumerate}

\section*{2. Two editions, one score}

Every piece is published as a critical edition and a performance edition.

\begin{enumerate}

\tightlist
\item
  \textit{The score is the same in both.} Note values, barring, voice names and editorial signs do not change. The performance edition simplifies the paratext. No reduced values, no modern barring, no piano reduction, no stripping of editorial signs.
\item
  \textit{Pitch is the one exception.} The critical edition stays at written pitch. The performance edition may be transposed when the written pitch does not suit a modern SATB choir, and its clefs follow the new ranges (an Altus in the octave treble clef becomes a treble clef). The masthead says so in one line, and a transposed edition drops the incipits, which would show the source's pitch.

  \begin{itemize}
  \tightlist
  \item
    \emph{High clefs (chiavette)}: transpose by the period's own interval, argued from period theory in the critical edition. \emph{Vox in Rama} goes a fifth down.
  \item
    \emph{Low clefs (chiavi naturali)}: no period rule asks for upward transposition; written pitch fixed only the relation between the voices, and the pitch was set by the director or the organ. Choose the interval from the voice ranges, by the time each voice spends in its comfortable range, not by its extremes alone. \emph{Już się zmierzka} goes a minor third up.
  \end{itemize}
\item
  \textit{Editorial suggestions} that go beyond the source (an octave drop at a final, say) appear only in the performance edition, as small notes in round brackets, explained in one sentence.
\item
  \textit{Final editions contain no working material.} No scripts, statistics, ‘consulted’ remarks, notes to self or status flags. The research record lives in this repository (issues and \texttt{notes/}).
\item
  \textit{Gaps.} Name only gaps in the primary sources (a lost print, an unseen facsimile, an uncollated witness), in one sentence in the Sources section. Unseen secondary literature is not listed; it goes to the issue tracker.
\end{enumerate}

\section*{3. Sources and sigla}

\begin{enumerate}

\tightlist
\item
  \textit{Sigla} of sources and editions are set in a thin ring (one letter) or rounded frame (several), so that a source A is never read as the word ‘a’ and never confused with the voice A. Voice abbreviations are plain small capitals.
\item
  \textit{Base text} is the earliest complete witness that can be read. Where that is a facsimile of a lost source, the facsimile is the base and the lost source keeps its own siglum.
\item
  \textit{Choosing sigla.} Short capitals. Use a letter for the source's own name or shelfmark (T = Tenor primi chori; K = Kraków BJ 1619). Where older literature numbers sources differently, follow the shelfmark and say why in one sentence.
\item
  \textit{Collated editions} get sigla too (PWM, K, M, R, N), listed separately under ‘Editions collated’.
\item
  \textit{Derivative editions.} Agreement in good readings proves nothing, since those go back to the source. The evidence is shared errors and shared innovations: readings, accidentals or underlay that two editions print and the source lacks. Editions that share them descend from one another or from a common parent, and count as one witness. Name the shared readings the argument rests on. (Grier 63--64, 77--79, 88)
\item
  \textit{Recordings} may be collated for accidentals and underlay ‘as heard’, with a siglum (P = Schola Gregoriana Pragensis, \emph{Plaude}). The edition's reading still rests on the source and the counterpoint, not on the recording.
\item
  \textit{Lost voices.} Where only some voices survive in the source, the title page and method state which voices come from the source and which from modern editions.
\item
  \textit{Earlier readings.} When an editor's change (Perz, say) is weighed against the source, each critical note ends with a verdict in small capitals: SOURCE RESTORED or {[}EDITOR{]} ACCEPTED.
\item
  \textit{Source descriptions.} Every source, and every copy of a print that was consulted, is named by city, library and shelfmark, in the library's own form; a lost source by its last recorded shelfmark. A printed title page is transcribed in full, in its own spelling, with its line ends marked ` \textbar{} '. (Grier 55--56, 212, 227--229; Caldwell 6--7)
\end{enumerate}

\section*{4. Pitch}

\begin{enumerate}

\tightlist
\item
  \textit{Print at written pitch.} The critical score keeps the source's final, signature and clef relations.
\item
  \textit{Pitch must be semi-reasonable for singers.} ‘All for historical pitch, but the pitch has to be semi-reasonable for singers to begin with.’ The performance notes give the written ranges and the transposition a modern choir will usually take; the notation does not change.
\item
  \textit{Unclefed sources.} Choose the pitch level by testing the counterpoint: the right level gives consonance on the strong beats; moving any one voice by a step gives dissonance throughout. State that pitch and signature are editorial.
\item
  \textit{Chiavette}: see 2.2. Evidence for the interval (fourth or fifth) comes from period theory (Banchieri, Picerli, Doni, Praetorius), the bass range test and the notation the transposition produces.
\item
  \textit{Octave of a voice.} Test with voice-crossing counts against neighbouring voices before accepting or emending an octave.
\end{enumerate}

\section*{5. Note values, barring, signs}

\begin{enumerate}

\tightlist
\item
  \textit{Unreduced values.} Note values are those of the source. Where the source is an organ tablature that halves its vocal model, double them.
\item
  \textit{Bar = breve} under cut-C. Bar numbers count breves, including both endings of a repeat.
\item
  \textit{Mensurstriche.} Bar lines run between the staves only, dashed, so that the score reads as parts first and as a timed score second. A note that crosses a bar line keeps its single value; the edition adds no ties. Section, repeat and final bar lines are solid. A bar line never passes through written text (a syllable, a punctuation mark, a stanza number); where it passes between the syllables of a word, the hyphen or extender leaves a gap for it.
\item
  \textit{Mensuration sign} as printed in the source. Where the source has none (a tablature, an unclefed manuscript), the edition may supply one, in square brackets ({[}¢{]}), and says in the method that it is editorial (Caldwell 50). The score must read without it: the light bar lines carry the time for those who want it.
\item
  \textit{Incipits} before each staff give the source's clef, sign and first note wherever a mensural source gives them. None where the source is a tablature, unclefed or lost.
\item
  \textit{Ambitus} of each voice after the clef.
\item
  \textit{Longs.} A long at a final is the period's way of saying what a modern editor says with a fermata: hold until all voices have finished. So:

  \begin{itemize}
  \tightlist
  \item
    Print the source's long at every final: the end of the piece, and the end of a section, a \emph{pars} or a repeated half where the source marks the break.
  \item
    Add no fermata over a long. Print a fermata (corona) only where the source has one.
  \item
    A long inside a phrase is a measured value, not a hold. Print it as such.
  \item
    A long that must end with the other voices but enters late is scaled to its real length. A maxima in one voice only may be given as a long, with a note.
  \end{itemize}
\item
  \textit{Rests.} In the performance edition every rest is cut at the bar lines and grouped by the bar, so that entries can be counted; notes are never cut (5.3), because a cut note implies a tie, while a cut rest changes nothing that sounds. A rest that fills its bar is centred in it (Gould 159). The critical edition keeps the rests of a primary source as written, with their groupings; voices taken from a modern edition are regularised as in the performance edition.
\item
  \textit{Note-value equivalence.} Where the edition's values are not the source's (a tablature's values doubled, 5.1), an equivalence stands above the first system of the critical score: the edition's note, an equals sign, the source's note. (Caldwell 23, 50)
\end{enumerate}

\section*{6. Clefs and voice names}

\begin{enumerate}

\tightlist
\item
  \textit{Voice names} are the source's (Cantus, Altus, Tenor, Bassus; Discantus, Medius, Tenor), in Latin, in small capitals, abbreviated after the first system. A name the source does not give is editorial and stands in square brackets at its first appearance in the critical edition ({[}Cantus{]}); the performance edition prints it plain. (Caldwell 52)
\item
  \textit{Modern clefs} are treble, octave treble and bass. The original clefs appear in the incipit and the Sources entry.
\item
  ‘Tenor’ names a function, not a range.
\end{enumerate}

\section*{7. Accidentals}

\begin{enumerate}

\tightlist
\item
  \textit{Source accidentals} stand before the note, full size, without brackets.
\item
  \textit{Editorial accidentals} (\emph{ficta}) stand above the note, small, without brackets, and apply to that note only.

  \begin{itemize}
  \tightlist
  \item
    \textit{Optional accidentals} stand above the note in square brackets, since editorial matter is square-bracketed: the note as printed is the edition's reading, the bracketed sign an alteration singers may make. Use them where the period rules leave a real choice; give the reason in the critical notes. Paired alterations (a double leading-tone cadence) are marked in both voices and the notes say ‘both or neither’. (Berger 1987)
  \end{itemize}
\item
  \textit{Scope and cautionaries.} A sign governs only its own note. Where an altered pitch returns unaltered in the same bar of the same voice, a cautionary sign stands before the returning note in round brackets, so that no singer carries the alteration on. A cautionary changes no reading and needs no note. (Caldwell 59--60; Grier 163--164; Gould 86)
\item
  \textit{Grounds.} Editorial accidentals rest on period rules: no \emph{mi contra fa} between voices; a sixth opening to an octave at a cadence is major; \emph{una nota super la}. Where the rules conflict, the critical note explains the choice. Parallel fifths count as grounds for a natural.
\item
  \textit{Consistency.} At a double leading-tone cadence, sharpen both voices or neither, never one.
\item
  \textit{Every editorial accidental must be justified.} If an accidental appears in a modern edition and not in the source, it is editorial and must be weighed, not inherited.
\item
  When modern-edition accidentals must stand for a lost voice, the method says so.
\end{enumerate}

\section*{8. Signs for editorial work}

\begin{longtable}[]{@{}
  >{\raggedright\arraybackslash}p{(\columnwidth - 2\tabcolsep) * \real{0.5000}}
  >{\raggedright\arraybackslash}p{(\columnwidth - 2\tabcolsep) * \real{0.5000}}@{}}
\toprule\noalign{}
\begin{minipage}[b]{\linewidth}\raggedright
What
\end{minipage} & \begin{minipage}[b]{\linewidth}\raggedright
Sign
\end{minipage} \\
\midrule\noalign{}
\endhead
\bottomrule\noalign{}
\endlastfoot
Accidental supplied & small, above the note \\
Accidental optional & small, in square brackets, above the note (the note keeps the edition's reading) \\
Cautionary accidental & in round brackets, before the note \\
Note supplied by the editor & small notehead \\
Note illegible or missing, supplied & notehead in square brackets \\
Source note divided to carry text & dashed tie \\
Ligature & horizontal bracket over the notes, kept where the underlay breaks it (10.16) \\
Coloration & corner brackets {\accfont ⌜} {\accfont ⌝} \\
Text the source gives by \emph{ij} & written out, in angle brackets ⟨ ⟩ \\
Text the source lacks & italics \\
Mensuration sign supplied & in square brackets, {[}¢{]} \\
Voice name supplied (critical edition) & in square brackets, {[}Cantus{]} \\
Cantus firmus entry & {[}c.f.{]} above the staff \\
Editorial suggestion (performance edition only) & small note in round brackets \\
Emendation & silent in the score; critical note opens ‘Emended.’ (13.6) \\
\end{longtable}

Square brackets mark editorial matter, an optional accidental included. Round brackets mark what changes no reading: a cautionary, and a suggestion in the performance edition (2.3). Each sign has one meaning, and the same meaning in every edition of the series; the build checks enforce it. A sign is judged in its context, not by what it might mean in another kind of score: italic text under a repeat reads plainly as a repeat to a singer, so a sign is not changed only because a modern reader might meet it elsewhere with another meaning. (Gould 494)

\section*{9. Repeats and structure}

\begin{enumerate}

\tightlist
\item
  \textit{Each repeated section or half starts and ends on its own system.}
\item
  Where the music changes inside a repeat, write the changed bars into both endings so the variation can be read and described.
\item
  Liturgical order (antiphon, verse, Gloria Patri, antiphon repeated) is shown with rubrics in the score.
\end{enumerate}

\section*{10. Underlay}

\textit{Firm rules} (kept without exception unless the source itself breaks them):

\begin{enumerate}

\tightlist
\item
  A new syllable only on a minim or longer, or on a semiminim straight after a dotted minim, and then the note after that semiminim takes a syllable too; never on a fusa or a dot. Prefer the longer note where the music allows. The semiminim after a dotted minim, and the syllable on the note after it: Lanfranco (1533) allows it ‘rarely’ and gives a syllable to the white note after it as well; Zarlino's rules V--VI (1558) permit it by implication and let the next note take a syllable too; Vicentino (1555) calls it no great fault (Harrán 1973a, 40--41; 1973b, 623). Stoquerus reports that the older composers often texted it and the newer ones, under ¢, did not (Rotola ed., 201, 221). The sources do it (the 1611 Tenor of \emph{Vox in Rama}). Under C rather than ¢, semiminims count as minims and take syllables freely (Stoquerus, 227, 243; Harrán 1973a, 40).
\item
  A syllable on the first note after a rest.
\item
  The last syllable of a phrase on its last (cadence) note.
\item
  No rest divides a word.
\item
  \textit{Elision.} None in Latin (Vicentino 1555). In Polish, the syllable count of the verse decides: elide only where the line needs it to keep its count, as the poet's own metre shows. (Vicentino requires elision in the vernacular; Polish syllabic verse rarely needs it.)
\end{enumerate}

\textit{Method.}

\begin{enumerate}

\setcounter{enumi}{5}
\tightlist
\item
  \textit{Words and music together.} Neither leads by rule. Map the cadences and place each line's last syllable on its arrival note. Where the music declaims, the words shape it: the stressed syllable or the key word takes the long note or the high point. Where the line sings, let it: a long run is right where the syllable can carry it, a stressed or key syllable or an open one. Vowel difficulty in a given range (Vicentino's register scheme, Harrán 1973b, 628) is not in itself a reason to change the underlay. Follow the momentum of the language. Light words and syllables (\emph{et}, \emph{in}, \emph{de}, the unstressed syllable before a stress) tend to fall off the beat inside moving figures and to lead into the stressed or open syllable on the beat; on long notes, where the choice is textural rather than declamatory, they may stand on the beat. A rising octave is taken inside one syllable: start the syllable on the lower note rather than changing on the upper one, preferably by drawing the next syllable back onto the lower note rather than stretching the previous one into the leap. A falling octave may carry a change on its landing note, or not. Where the lower note is a cadence note carrying a phrase's last syllable (10.3), that rule wins. After a run of short notes, change syllable before the run rather than on the white note straight after it, where the notes allow (Lanfranco 1533, rule VI, in Harrán 1973a, 41). The voices need not change syllable together. These are tendencies to weigh by ear, not rules to count: sing every line before settling it. This matters most where an edition underlays from scratch.
\item
  \textit{Wacław's own practice} (Kraków \emph{Lamentationes}, 1553): mostly syllabic, repeated notes take new syllables; one melisma before each cadence, on the stressed syllable of the last word; the final syllable on the cadence note; short repeats \emph{ij}, longer ones written out.
\item
  \textit{Accent} follows the penultimate rule as taught in Kraków (Sebastian z Felsztyna, 1518).
\item
  \textit{Repeated notes} take a new syllable. The one exception (Stoquerus, ch.~15, Rotola ed., 175--181): two notes on the same pitch whose first is a semiminim or shorter and follows a note or dot of its own value; then the second takes the syllable, or better neither does. Long free melismas may be replaced by repeating the clause.
\item
  \textit{Imitation.} A point of imitation carries the words of its head motif wherever it recurs.
\item
  \textit{Source guides.} Where the source prints text in blocks, use its spacing (gaps mark melismas) and the voice whose text fits note for note.
\item
  \textit{Departures} from the theorists that the source itself shows are kept and noted.
\item
  \textit{Lower voices} that are editorially texted may skip words to keep long notes, as long as their text makes sense alone. Divide notes only where text requires it, mark the division with a dashed tie, and avoid movement before the penultimate bar at a final cadence.
\item
  \textit{Strophic songs.} Underlay as many stanzas as the music needs. Where the underlay is the same in every stanza, stanza 1 goes under the notes and the rest follow the score. Where a later stanza needs different underlay (a different stress pattern, an extra syllable, a melisma that would fall on a weak syllable), that stanza goes under the notes too. Weigh how the piece is usually performed: if most performances sing one or two stanzas, those two matter most.
\item
  Theorists are cited by paraphrase; direct quotation is not needed.
\item
  \textit{Ligatures and underlay.} A ligature of the source keeps its bracket even where the underlay puts a new syllable inside it, since the ligature is evidence for the underlay and the articulation. The critical note says so. (Caldwell 60--61)
\end{enumerate}

\emph{Literature for this section}: Harrán, D., `New Light on the Question of Text Underlay Prior to Zarlino', \emph{Acta Musicologica} 45 (1973a), 24--56; Harrán, `Vicentino and His Rules of Text Underlay', \emph{Musical Quarterly} 59 (1973b), 620--632; Harrán, `How to ‘Lay’ the ‘Lay’: New Thoughts on Text Underlay', \emph{Musica Disciplina} 51 (1997), 231--262; Towne, G., `A Systematic Formulation of Sixteenth-Century Text Underlay Rules', \emph{Musica Disciplina} 44 (1990), 255--287, and 45 (1991), 143--168; Stoquerus, Caspar, \emph{De musica verbali libri duo}, ed.~and trans. A. C. Rotola (Lincoln, 1988).

\section*{11. Text}

\begin{enumerate}

\tightlist
\item
  \textit{Latin spelling} is the source's. Minor variants that are likely errors (a letter dropped, doubled or misread) are normalised, and the critical notes give the source's form. Spellings that are more idiosyncratic, consistent in the source or typical of its time and place (\emph{confidencium}, \emph{misericordie}), are kept: they are part of the source's character. Typography is modernised silently: \emph{et} for \&, \emph{u} and \emph{v} distinguished, abbreviations expanded. Where the source has no text, or only an incipit, text and spelling follow the liturgical book (\emph{exspectatione}, \emph{Iudaeorum}). Punctuation from the relevant liturgical book (Missal 1570 for \emph{Vox}). (Caldwell 63)
\item
  \textit{Old Polish}: modern spelling and punctuation, old sounds and forms kept (\emph{zmierzka}, \emph{naszem}, \emph{wszytki}, \emph{swoję}, \emph{jenż}). Do not modernise to \emph{zmierzcha}, \emph{naszym}, \emph{On}.
\item
  \textit{Diplomatic text} in the critical edition keeps the source's spelling and line division (ſ kept, abbreviations unexpanded).
\item
  \textit{Translation} is literal, line by line. Archaic English only where the source language is liturgical verse.
\item
  \textit{Pronunciation} lives in one series guide (\texttt{guides/pronunciation}), not in each edition. An edition notes only what is peculiar to its text (a word that needs a gloss, an option the text invites) and points to the guide. Defaults: Polish Latin for Latin texts; modern pronunciation of the old forms for Polish. Tables, never paragraphs.
\item
  \textit{Lemmata} quote the source's spelling, including a spelling the edited text normalises (11.1).
\item
  \textit{Translations of quotations.} Every quotation in a language other than English, in either edition, carries an English translation inline: after the original, in parentheses and quotation marks, `\ldots{}' (`\ldots{}'). A block quotation is followed by its translation inside the same block. Titles of works are translated at their first mention in the prose where the meaning matters to the argument; not in source lists or the literature.
\item
  \textit{Latin syllable division} follows one system throughout. A consonant group that can begin a Latin word goes with the following syllable (\emph{o-mnis}, \emph{San-cto}, \emph{co-gno-vi-sti}, \emph{ex-spe-cta-ti-o-ne}); other groups divide (\emph{sem-per}, \emph{an-ge-lum}); a compound divides at its prefix (\emph{ex-spe-}, \emph{in-spi-ce}). Singers carry the consonant over to the next syllable, so that the vowel before it turns quickly, and the hyphen shows where they sing it. Polish follows the modern dictionary. (Caldwell 42; Gould 445--446)
\item
  \textit{Text written \emph{ij}.} Words the source gives by \emph{ij} or another repeat sign are written out under the notes, in angle brackets ⟨ ⟩. Italics are kept for text the source lacks entirely. (Caldwell 63, 103)
\end{enumerate}

\section*{12. Chant}

\begin{enumerate}

\tightlist
\item
  Base text and melody from the oldest notated witness; neume grouping follows it.
\item
  Square notation, unmeasured. Mensural or modern rhythmic readings are reported, not adopted.
\item
  Bar lines are editorial and mark verse structure. \emph{Mora vocis} only at stanza ends.
\item
  No B-flat unless the base witness has it.
\item
  Faint or invisible notes: say so in a note; supply from parallels and earlier transcriptions.
\item
  Later hands are described, not normalised away.
\end{enumerate}

\section*{13. Paratext}

\textit{Critical edition}: Title page, colophon and contents, Introduction (the piece, composer, poet, context), Sources, Editorial method, Text and translation, {[}Facsimile{]}, Score, Critical notes, Literature.

\textit{Performance edition}: Masthead (with a pitch line if transposed), Score, stanzas after the score, then Text and translation and For rehearsal (the song, the words, stanzas, pitch and voices, tempo, the score's signs), colophon line. No context before the music: setting the scene is the director's job, and the rehearsal notes give them what they need for it.

\begin{enumerate}

\tightlist
\item
  Context is welcome in both: ‘enough paratext for both performers and critical readers’. In the performance edition it is limited to what changes how one sings.
\item
  Distinguish verified citations from the editor's own readings. Write ‘my own reading’ where no published study makes the link.
\item
  Verify every claim before printing it; cut what cannot be verified.
\item
  \textit{Do not restate the series rules in an edition.} The rules live in one place, the series guide \emph{Editorial principles} (this document, also kept in the repository as docs/editorial-principles.md). An edition's method section says only what is particular to the piece: its sources, the choices the rules leave open, and every departure from a rule, with the reason. Where a rule is meant, name it and cite it by number (‘principles 10.1’). Restating rules makes the editions repetitive and lets them drift out of step with each other and with this file.
\item
  \textit{The performance edition argues from the music, not from names.} It gives the reason for a choice in general terms (‘as an organ would have done’), never by pointing to what a particular ensemble, recording or modern edition does. Specific citations, if they matter, belong in the critical edition.
\item
  \textit{Emendations.} A critical note that opens ‘Emended.’ gives the source's reading and says how the error probably arose: a clef or octave slip, an eye that skipped to a neighbouring part or line, a misread tablature letter, a stem left out. Where the evidence does not show the cause, the note says that it is unclear. (Grier 72, 99--100; Caldwell 5)
\item
  \textit{Tense.} Critical notes, Sources and method describe an extant source in the present tense (`T prints e′') and a lost or destroyed one in the past (`A had a maxima'). (Caldwell 9)
\end{enumerate}

\section*{14. Names, attribution, rights}

\begin{enumerate}

\tightlist
\item
  \textit{Composers by their vernacular names only} (Piotr z Grudziądza, Wacław z Szamotuł, Mikołaj Zieleński). Latin forms appear only where the prose discusses them, such as a signature in a source. An anonymous work is credited to \emph{Anonymous}, with its date in the composer's dates slot.
\item
  \textit{Title and subtitle.}

  \begin{itemize}
  \tightlist
  \item
    \emph{Title}: the work's title in its main source (a heading, a rubric, a table of contents), in the source's language. If the source gives none, the incipit.
  \item
    \emph{Subtitle}: if the title is not the incipit, the incipit (\emph{Modlitwa gdy dziatki spać idą} / \emph{Już się zmierzka}). If the title is the incipit, the work's liturgical or functional designation in the source's language (\emph{Vox in Rama} / \emph{Communio in festo SS. Innocentium}). If neither applies, no subtitle.
  \item
    Alternative titles go in the introduction. The collection or print a piece belongs to goes in the source line. Scoring goes in the forces line (‘for four voices’), never in the subtitle.
  \end{itemize}
\item
  \textit{Editor}: Mikołaj Derduń, with ń.
\item
  \textit{Series}: Polish Early Music, numbered in order of first edition.
\item
  \textit{Licence}: editorial text, scores and sources under CC BY 4.0; build code under MIT. The music is public domain. Facsimile images only where their holder's terms allow (Leipzig and Polona images are public domain).
\item
  \textit{Edition numbering}: ‘First edition, 2026’. Corrections that change readings make a new edition; typographic fixes do not.
\end{enumerate}
